\documentclass{article}
\usepackage{amsmath,amssymb,amsthm,mathtools}
\newtheorem{theorem}{Theorem}
\newtheorem{lemma}{Lemma}
\newtheorem{proposition}{Proposition}
\newtheorem{corollary}{Corollary}
\newtheorem{definition}{Definition}
\newtheorem{example}{Example}
\title{ShadowBench geometry/L2/geo\_gen\_L2\_003}
\begin{document}
\maketitle
% Problem id: geometry/L2/geo_gen_L2_003
% Companion instructions: docs/instructions.md
% Candidate skeletons: docs/skeletons/
% Lean target: ShadowBench/Source/Main.lean

\begin{theorem}[erdos_mordell_inequality]\label{thm:erdos_mordell}
In Euclidean geometry, the Erdős–Mordell inequality states that for any triangle $ABC$ and point $P$ inside $ABC$, the sum of the distances from $P$ to the sides is less than or equal to half of the sum of the distances from $P$ to the vertices.
Let $PL, PM, PN$ be the perpendiculars from $P$ to the sides $BC, CA, AB$ respectively. Then:
\[
PA + PB + PC \ge 2(PL + PM + PN)
\]
\end{theorem}
\end{document}
